\chapter{Parallelism and HPC Deployment}
\label{ch:hpc}

\qv{} parallelises at three nested levels: OpenMP threads inside one solver
call, processes across reads on one node, and MPI/SLURM across nodes. They
compose --- but must be composed deliberately to avoid oversubscription.

\section{OpenMP: inside a run}

Built in by default (\code{QANNEAL\_ENABLE\_OPENMP=ON}). Parallel regions:

\begin{center}
\small
\begin{tabular}{@{}lp{8.6cm}@{}}
\toprule
\textbf{Class} & \textbf{Parallelised over}\\
\midrule
\code{ReplicaAnnealer} & replicas\\
\code{SQAAnnealer} & replicas (each owning its slices), when no sweep-level
observer is attached\\
\code{SQAParallelTemperingAnnealer} & ladder replicas\\
\bottomrule
\end{tabular}
\end{center}

Thread count is controlled at runtime:

\begin{bashcode}
export OMP_NUM_THREADS=8
\end{bashcode}

Each replica owns a private RNG stream, local-field cache and cluster
buffers, so results are independent of the thread count; runs are
reproducible for a fixed \code{seed} and fixed replica count.

\begin{warnbox}[Oversubscription]
Threads only help while \code{OMP\_NUM\_THREADS} $\le$ the number of
parallel units (replicas). And when you add process-level parallelism
(below), \emph{multiply} the two: 8 worker processes $\times$ 8 OpenMP
threads = 64 threads on what may be a 16-core node --- a large slowdown, not
a speedup. Standard practice on shared nodes: set
\code{OMP\_NUM\_THREADS=1} for many-process runs, or partition cores
explicitly (e.g.\ 4 workers $\times$ 4 threads on 16 cores).
\end{warnbox}

\section{Process-level: the HPC launcher}

\code{examples/python/hpc\_sqa\_launcher.py} distributes \emph{reads}
(independent restarts) over workers --- the ideal parallelisation axis, since
reads never communicate:

\begin{bashcode}
# Single node, multiprocessing (no extra dependencies)
python examples/python/hpc_sqa_launcher.py \
  --n 5000 --method sqapt --reads 64 --workers 8

# Multi-node MPI (requires mpi4py)
mpirun -n 32 python examples/python/hpc_sqa_launcher.py \
  --n 50000 --method sqa --reads 8 --mpi

# Built-in scaling study
python examples/python/hpc_sqa_launcher.py --scaling --method sqapt --mode fast
\end{bashcode}

Each worker runs a full solver call with a distinct seed and writes a JSON
result; the launcher merges them and reports the global best.

\section{SLURM job arrays}

The most portable route on clusters --- no \code{mpi4py}, no long-running
master, natural fault tolerance (a failed task costs one read batch):

\begin{bashcode}
sbatch scripts/slurm/run_sqa_array.sh
# afterwards, merge all task outputs:
python examples/python/merge_array_results.py results/run_*.json
\end{bashcode}

Adapt the array script's template variables (problem size, method, reads per
task, walltime) to your site; give each array task a distinct
\code{--seed} derived from \code{\$SLURM\_ARRAY\_TASK\_ID} so reads are
independent.

\section{C++ MPI}

For fully native distributed annealing:

\begin{bashcode}
cmake -S . -B build -DQANNEAL_ENABLE_MPI=ON
cmake --build build -j
mpirun -n 8 build/qanneal_mpi_example
\end{bashcode}

The MPI layer (\code{MPIContext}, \code{run\_replica\_anneal()} in
\code{include/qanneal/mpi/}) runs one replica group per rank and reduces the
best result to rank 0. See \code{examples/mpi/sa\_mpi.cpp} for the pattern.

\section{Choosing the axis of parallelism}

\begin{center}
\small
\begin{tabular}{@{}lp{5.2cm}p{5.6cm}@{}}
\toprule
\textbf{Level} & \textbf{Use when} & \textbf{Cost model}\\
\midrule
OpenMP (replicas) & one hard instance, want a better $\chiB$ estimate or PT
ladder & shared memory; near-linear up to replica count\\
Processes (reads) & success probability per read $<1$; want
$1-(1-p)^R$ aggregation & embarrassingly parallel; linear\\
Nodes (MPI/SLURM) & instance sweeps, many problems $\times$ many reads &
linear; prefer job arrays for robustness\\
\bottomrule
\end{tabular}
\end{center}

A practical default for one hard instance on one node: \code{replicas=8}
with \code{OMP\_NUM\_THREADS=8}, and \code{reads} spread over remaining
cores via the launcher.
